\naslov{Statistično sklepanje z nadzorovanim tveganjem}

Opazimo $X$, zanima nas $Y$.
Želimo zatrditi, da z visoko verjetnostjo velja $Y_{\text{min}} < Y <
Y_{\text{max}}$, kjer sta $Y_{\text{min}}$ in $Y_{\text{max}}$ opazljivi.
Intervalu med njima pravimo \pojem{napovedni interval}, če pa je $Y = y =
g(\theta)$ determinističen, pa \pojem{interval zaupanja}.

Splošneje lahko izjavimo $Y \in C_X$, kjer je $\{C_X\}_X$ nek nabor množic.
Naš cilj je nadzorovati verjetnost zmote $p_\theta = P_\theta(Y \notin C_X)$, za
kar vnaprej izberemo \pojem{stopnjo tveganja} $\alpha$, tipično $0.05$.
To lahko naredimo na več načinov:
\begin{itemize}
\item eksaktna stopnja tveganja: $p_\theta = \alpha$ za vsak $\theta$,
\item nominalna stopnja tveganja: $p_\theta$ je čim bližje $\alpha$,
\item asimptotično eksaktna stopnja tveganja: opazimo čim več vrednosti;
  zahtevamo
  \[
	\lim_{n \to \infty} p_\theta = \alpha,
  \]
\item konzervativna stopnja tveganja: $p_\theta \le \alpha$ za vsak $\theta$.
\end{itemize}
Nasprotni verjetnosti $1 - \alpha$ pravimo \pojem{stopnja zaupanja}.

\vprasanje{Na katere načine lahko izberemo stopnjo zaupanja?}

Za iskanje mej definiramo \pojem{mero razhajanja} $\rho(x,y)$, ki meri neskladje
med $x$ in $y$.
Če imamo srečo, je $\rho$ \pojem{pivotna funkcija}, kar pomeni, da je
porazdelitev slučajne spremenljivke $\rho(X,Y)$ neodvisna od $\theta$.
Tedaj postavimo $C = \{(x,y) \such \rho(x,y) < c_\alpha\}$, kjer $c_\alpha$
izračunamo tako, da je verjetnost zmote enaka $\alpha$.

\begin{izrek}
  Če definiramo
  \[
	\widehat{\text{RMSE}}
	= \sqrt{\inv{n} \grad g(\hat{\theta})^T
	  \operatorname{FI}_1^{-1}(\hat{\theta}) \grad g(\hat{\theta})},
  \]
  kjer je $\hat{\theta}$ cenilka po MNV, pod podobnimi pogoji kot za prejšnji
  rezultat velja
  \[
	\frac{\hat{y} - y}{\widehat{\text{RMSE}}}
	\xrightarrow[n \to \infty]{d} N(0,1).
  \]
\end{izrek}

% LocalWords:  MNV
